\begin{abstract}
\textbf{Purpose} JIGSAWS sanctions two cross-validation protocols for surgical skill assessment: one holds out a repetition while every surgeon remains in training, the other holds out a surgeon entirely. We ask what a model evaluated under the first measures, and whether reported confidence would support deferring uncertain trials to a human rater.

\textbf{Methods} Two prediction paths were evaluated on identical information across 103 trials, 8 surgeons and 3 tasks: 16 kinematic features with ridge and multinomial logistic regression, and a zero-shot commercial decision model given the same features rendered as text. Both protocols were run. Errors are reported against a feature-free reference predicting each held-out trial with its training-fold mean rating; intervals resample whole surgeons. Analysis rules were fixed in advance; post-hoc additions are labelled.

\textbf{Results} Changing only the protocol removed at least a third of skill-classification accuracy (paired gap 0.49-0.67, intervals excluding zero). A post-hoc diagnostic found the same features identified which of eight surgeons performed a held-out trial with accuracy 0.83-0.93 against chance near 0.13, separating identity confounding from underfitting. On unseen surgeons one task of three beat the reference, and there a single-feature regression on trial duration matched both models. Both paths were overconfident (confidence 0.755 against accuracy 0.379; 0.696 against 0.301), scoring worse than a constant forecast, and neither identified deferrable trials.

\textbf{Conclusion} Accuracy reported under the permissive protocol is consistent with surgeon recognition rather than skill measurement. With eight surgeons these are single-dataset demonstrations rather than established generalities. We recommend reporting held-out-surgeon results against an uninformative reference, with a subject-identifiability diagnostic.

\textbf{Keywords} Surgical skill assessment; JIGSAWS; identity confounding; cross-validation; calibration; selective prediction

\end{abstract}
